\documentclass[10pt]{beamer}

\usetheme[%
    progressbar=frametitle,
    block=fill,
    numbering=fraction,
    footer=crumbs,
    sectionpage=numbered,
    subsectionpage=none,
    titleformattitle=smallcaps,
    titleformatsubtitle=smallcaps,
    %%% new options for UR2 theme:
    maincolor=red,
]{metropolis-ur2}

\usepackage{hyperref}
\hypersetup{
    colorlinks=true,
    linkcolor=.,
    filecolor=ur2Turquoise,
    urlcolor=ur2Bleu,
    pdftitle={M1 MAS - Python - Généralités},
    pdfpagemode=FullScreen,
}

% \usepackage{biblatex}
% \addbibresource{example.bib}

\usepackage{appendixnumberbeamer}
\usepackage{adjustbox}
\usepackage{booktabs}
\usepackage{fontawesome5}
\usepackage{graphicx}
\usepackage{array}
\usepackage{tabularx}
\usepackage{myminted}
% \usepackage[cache=false]{minted}

\title[M1 MAS -- Python -- Généralités]{CM 0 - Généralités}
\subtitle{Programmation Python -- Master 1 MAS}
\author{Romain Tavenard}
\date{2026}
\institute{%
\hypersetup{urlcolor=.}
\makebox[2.2ex][c]{\faEnvelope}\enspace\href{mailto:romain.tavenard@univ-rennes2.fr}{\texttt{romain.tavenard@univ-rennes2.fr}}\\%
% \makebox[2.2ex][c]{\faHome}\enspace\url{https://rtavenar.github.io/}%
}

\begin{document}

\maketitle

\begin{frame}{Format du cours}  
  \begin{itemize}
    \item 12 semaines
    \begin{itemize}
      \item 1h CM, 1h TD
    \end{itemize}
    \item Travail personnel à prévoir
    \begin{itemize}
      \item 2h par semaine en moyenne (donc + pour certain(e)s)
      \item Projets
      \item Préparation de cours + QCM
    \end{itemize}
    \item Évaluation
    \begin{itemize}
      \item 2 contrôles continus (1 papier, 1 sur machine)
      \item 1-2 points bonus pour les étudiants qui auront effectué tous les QCM en temps et en heure (barême à préciser)
      \item Méthode de calcul de la note finale encore inconnue
    \end{itemize}
  \end{itemize}
\end{frame}

\begin{frame}{Contenu du cours}  
  \begin{itemize}
    \item 4 grandes parties
    \begin{enumerate}
      \item Bases du langage
      \item Structures de données (listes, dictionnaires, \emph{etc.})
      \item Accès aux données (fichiers, requêtes HTTP)
      \item Programmation Orientée Objet
    \end{enumerate}
  \end{itemize}
\end{frame}

\begin{frame}{Organisation (1/2)}  
  \begin{itemize}
    \item À chaque nouvelle partie
    \begin{itemize}
      \item Les TD + Projet de la partie sont mis en ligne
      \item Pas de projet pour la partie 1
      \item Autres projets non notés (mais avec \emph{feedback})
    \end{itemize}
    \item Avant chaque CM
    \begin{itemize}
      \item Un ou plusieurs chapitres de cours à travailler (polycopié)
      \item Un QCM CURSUS à remplir
    \end{itemize}
    \item Après chaque TD
    \begin{itemize}
      \item Finir le TD par soi-même
    \end{itemize}
  \end{itemize}
\end{frame}

\begin{frame}{Organisation (2/2)}  
  \begin{itemize}
    \item En séance de CM
    \begin{itemize}
      \item Focus sur quelques points saillants du cours
      \item Venir avec ses questions sur le cours
      \item Venir avec ses questions sur le projet
    \end{itemize}
    \item En séance de TD, par ordre de priorité :
    \begin{enumerate}
      \item Le TD du jour
      \item Le(s) TD(s) suivant(s) de la partie en cours
      \item Le projet de la partie en cours
    \end{enumerate}
  \end{itemize}
\end{frame}

\begin{frame}{Dualité TD / Projet}  
  \begin{itemize}
    \item Pourquoi TD \alert{et} projet ?
    \begin{itemize}
      \item Hétérogénéité de niveau en Python (début d'année)
      \item Apprendre un langage de programmation implique de la répétition
      \item Python sera un (LE ?) langage central de vos études en Master
    \end{itemize}
    \item Principe de base :
    \begin{itemize}
        \item Je sais faire les TD sans problème $\Rightarrow$ je vise 10/20
        \item Je sais faire les TD et les projets sans problème $\Rightarrow$ je vise 20/20
    \end{itemize}  
  \end{itemize}
\end{frame}

\begin{frame}{Quid des LLM ?}  
  \begin{itemize}
    \item A-t-on le droit aux LLM pour ce cours ? \alert{NON}
    \item Pourquoi ?
    \begin{itemize}
      \item Cours d'introduction à Python
      \item Coder avec LLM ne peut se faire que si on connaît les bases du langage (but de ce cours)
    \end{itemize}  
    \item Usage des LLM pour les révisions
    \begin{itemize}
      \item Possible pour se faire ré-expliquer un concept
      \item Possible pour se faire générer des exercices supplémentaires + demander une correction \alert{UNE FOIS QUE VOUS AVEZ FAIT PAR VOUS MÊME}
    \end{itemize}  
  \end{itemize}
\end{frame}

\end{document}
